\textbf{COLUNI 2019 -- Questão 33}\par\medskip A cidade de Brasília, um caso especial na rede urbana, foi planejada com a intenção de ser a capital da nação com seus tons de modernidade, funcionalidade e harmonia. A cidade foi projetada com racionalidade técnica, como um projeto geopolítico de ocupação do oeste e afastamento do poder político do litoral brasileiro.

\par

É CORRETO afirmar que a cidade de Brasília:\begin{enumerate}
\item eliminou a tradicional estruturação espacial das cidades como o contraste entre áreas reservadas à classe
      média e às elites e as reservadas às camadas populares.
\item fracassou como polo de atração populacional para o oeste, pois serviu apenas para uma camada de
      burocratas políticos habitarem.
\item concretizou-se como uma cidade onde problemas de trânsito não são registrados, em função das grandes
      vias, em especial o metrô previsto no planejamento.
\item consolidou-se como uma metrópole nacional, conectada ao restante do país pelo fato de ser a sede do
      governo federal.\end{enumerate}
